\section{Тренировка системы детекции объектов}

\subsection{Подготовка данных для обучения}

Подбор набора данных для тренировки любой нейросети требует наличие достаточного количества разнообразных данных.
При выборе следует рассматривать специальные наборы данных которые разрабатываются специально под определенные сферы. 
Для задачи детекции в автономном вождении сформулируем ряд критериев:

\begin{itemize}
    \item Релевантность классов --- наличие меток для транспортных средств (легковые, грузовики, автобусы), пешеходов, велосипедистов и тд;
    \item Вариативность --- разнообразие условий (освещение, погода), плотность трафика, географические особенности;
    \item Качество разметки --- точность ограничивающих рамок, согласованность аннотаций.
\end{itemize}

Для автономного транспорта существует не так много наборов данных, которые в полной мере удовлетворяют данным критериям. 
В качестве подходящего набора выбран набор KITTI, который удовлетворяет всем требованиям: содержит более 15 тысяч изображений
с размеченными объектами на все основные классы, которые необходимо распознавать (автомобили, велосипеды, мотоциклы), а также
имеет довольно редкие категории (трамваи, микроавтобусы).

Все данные KITTI собраны в реальных дорожных условиях с использованием лидаров, радаров и камер высокого разрешения. Популярность данного
набора в научном сообществе также указывает на его качество.